\chapter{The Standard Model}
\label{chap:the-standard-model}

% Closed question: What remains invariant under admissible relabeling or
% recompilation?

\begin{chapteressay}

A speedometer does not care which way the car is pointing. Northbound,
eastbound, or turning through a roundabout, the speed reading is a magnitude.
Rotate the coordinate axes and the speed should not change.

The odometer is different. It does care about the direction of time. Drive a
mile forward and a mile backward, and the odometer reads two miles, not zero.
Some transformations preserve the record. Others do not. A physical theory has
to know the difference.

This is the doorway into the Standard Model: not particles as a zoo of names,
but invariance as a measurement discipline. A relabeling is admissible when it
changes the description without changing the physical record. A relabeling is
not admissible when it changes what the apparatus can measure. Gauge symmetry
is the precise version of that rule. It tells the theory which local changes of
description must leave the physics unchanged and which fields are required to
make that invariance possible.

The second example is compilation. A program can be rewritten, optimized, and
compiled. If the resulting object has the same observable behavior, the
rewriting was admissible. If the output changes, the rewrite was not merely a
new description; it changed the measurement. The compiler is a laboratory for
invariance.

The code anchors Chapter 10 in exactly that way. \lean{LOGICAL} holds an
\lean{EKG} calibration. \lean{ComputerProgram.le} compares only the halted
boolean surface that can be calibrated. \lean{ElaborationProcess} turns
\lean{load} into \lean{transform} and \lean{transform} into \lean{boolean}.
\lean{HALTED} recognizes the boolean form as the only halted case.
\lean{MEASURED} and \lean{COMPILED} then carry the result forward as a measured
compiler record.

The closed question is:

\begin{center}
\emph{What remains invariant under admissible relabeling or recompilation?}
\end{center}

The answer is the physical content of the record. Noether's theorem makes
symmetry into conservation. Electromagnetic gauge invariance makes charge
conservation structural. Spin measurements reveal representations rather than
continuous hidden directions. Strong and weak interactions show that not every
intuitive symmetry survives measurement.

\NB{The current Lean code does not formalize U(1), SU(2), SU(3), Dirac
operators, Higgs fields, or the Standard Model Lagrangian. This chapter uses
Standard Model examples to explain the physical gauge concept, while the code
anchors the measurement and compiler interfaces that such a formalization would
eventually need.}

Chapter 10 therefore asks what survives when names change. General relativity
made frame comparison physical. The Standard Model makes internal relabeling
physical.

\end{chapteressay}

% Phenomenon arcs use: 1/3 setup -> phenomenon box -> 2/3 structural interpretation.

\section{Noether's Theorem}
\label{se:noethers-theorem}

Noether's theorem, proved by Emmy Noether in 1918, is the hinge between
symmetry and conservation. The statement: every continuous symmetry of an
action functional corresponds to a conserved current. The theorem applies
to all classical field theories and, in modified forms, to quantum field
theories as well.

The structural content is precise. Consider a classical system described by
an action $S[\phi] = \int L(\phi, \partial_\mu \phi) \, d^4x$, where $\phi$
is a field (scalar, vector, spinor, etc.) and $L$ is the Lagrangian density.
Suppose the action is invariant under a continuous transformation
$\phi \to \phi + \epsilon \delta\phi$, parameterized by a small constant
$\epsilon$. Noether's theorem then states that there exists a conserved
current $j^\mu$ satisfying $\partial_\mu j^\mu = 0$, and a corresponding
conserved charge $Q = \int j^0 \, d^3x$ that is constant in time.

The derivation is brief but worth retracing. Vary the field by $\epsilon
\delta\phi$ and compute the change in the Lagrangian density:
\[
\delta L = \frac{\partial L}{\partial \phi} \delta\phi + \frac{\partial L}{\partial(\partial_\mu \phi)} \partial_\mu \delta\phi.
\]
The Euler-Lagrange equations are $\partial_\mu(\partial L/\partial(\partial_\mu \phi)) = \partial L/\partial \phi$, so on shell:
\[
\delta L = \partial_\mu\left(\frac{\partial L}{\partial(\partial_\mu \phi)} \delta\phi\right).
\]
If the transformation is a symmetry of the action ($\delta L = 0$, or
$\delta L = \partial_\mu K^\mu$ for some current $K^\mu$ that integrates
to zero at the boundary), then the conserved current is
\[
j^\mu = \frac{\partial L}{\partial(\partial_\mu \phi)} \delta\phi - K^\mu.
\]
The 4-divergence $\partial_\mu j^\mu = 0$ is the conservation law.

The applications are foundational. Time-translation symmetry ($t \to t +
\epsilon$) gives energy conservation: $E = \int T^{00} d^3x$ is constant.
Space-translation symmetry ($\mathbf{x} \to \mathbf{x} + \boldsymbol{\epsilon}$)
gives momentum conservation: $\mathbf{P} = \int T^{0i} d^3x$ is constant.
Rotational symmetry gives angular momentum conservation. Each spacetime
symmetry produces a specific conserved tensor component of the
stress-energy-momentum tensor $T^{\mu\nu}$.

For internal symmetries, the structure is parallel. The most important
example is the U(1) gauge symmetry of electromagnetism. The free
electromagnetic Lagrangian is $L = -\frac{1}{4} F_{\mu\nu} F^{\mu\nu}$,
where $F_{\mu\nu} = \partial_\mu A_\nu - \partial_\nu A_\mu$ is the
electromagnetic field strength. This Lagrangian is invariant under the
gauge transformation $A_\mu \to A_\mu + \partial_\mu \chi$ for any
function $\chi$. When coupled to a charged matter field $\psi$ via the
covariant derivative $D_\mu = \partial_\mu - ieA_\mu$, the full theory
becomes invariant under a local U(1) phase rotation $\psi \to e^{ie\chi}
\psi$. Noether's theorem applied to this local symmetry gives the
conservation of electric charge: $\partial_\mu j^\mu = 0$ where $j^\mu$
is the electric current density.

The structural lesson is the chapter's central claim: the conservation
laws of physics are not contingent facts about nature. They are the
algebraic consequences of the symmetries of the action. Energy is
conserved because the laws of physics do not depend on the choice of
zero of time. Momentum is conserved because the laws do not depend on
the choice of origin in space. Electric charge is conserved because the
electromagnetic theory is invariant under local phase rotations of the
charged fields.

\begin{phenom}{Noether Effect}
\label{ph:noether-effect}

\PhStatement A continuous admissible relabeling of the description produces a
conserved physical record.

\PhOrigin Noether's theorem proved that differentiable symmetries of the action
correspond to conservation laws.

\PhObservation When the action is unchanged by time translation, energy is
conserved; when it is unchanged by phase relabeling in electromagnetism, charge
is conserved.

\PhConstraint The theory may not treat symmetry as cosmetic if the symmetry
forces a conserved measurement.

\PhConsequence Invariance is operational: what survives admissible relabeling
is carried forward as a conserved ledger entry.

\end{phenom}

This is why the chapter belongs after transport failure. Once transport can
twist labels, the theory must say which labels were only coordinates and which
differences were real. Noether gives the test: the admissible relabeling leaves
behind a conservation law.

\section{The Photoelectric Effect}
\label{se:the-photoelectric-effect}

The photoelectric effect forced a change in what light could be. Below a
threshold frequency, increasing intensity does not eject electrons. Above the
threshold, electrons emerge with kinetic energy

\[
K_{\max} = hf - \phi.
\label{eq:chapter10-photoelectric}
\]

The result says that energy exchange is quantized. The electron does not
accumulate arbitrarily weak light until it escapes. It absorbs a quantum whose
frequency must be high enough to pay the work function.

\begin{phenom}{Photoelectric Effect}
\label{ph:photoelectric-effect}

\PhStatement A measurement can reveal that a conserved exchange occurs in
quanta rather than by continuous accumulation.

\PhOrigin Einstein explained photoelectric emission by treating light energy as
localized quanta proportional to frequency.

\PhObservation Electron emission depends on frequency threshold, while
intensity changes the number of emitted electrons rather than rescuing
below-threshold light.

\PhConstraint The apparatus may not model electromagnetic energy transfer as
arbitrarily divisible when the emission record has a threshold quantum.

\PhConsequence Quantized exchange becomes part of the invariant measurement
content later organized by quantum electrodynamics.

\end{phenom}

In gauge language, the photoelectric effect is not the whole theory. It is a
door. It says that the invariant record of electromagnetic interaction is not a
smooth smear of energy. It arrives in countable exchanges.

\section{Spin-1/2 and the Stern--Gerlach Experiment}
\label{se:spin-half-and-stern-gerlach}

The Stern-Gerlach experiment sends silver atoms through an inhomogeneous
magnetic field. A classical magnetic moment might be expected to smear into a
continuous distribution. Instead, the beam splits into discrete outcomes.

For spin one-half, measurement along an axis gives two values, often written
as

\[
m_s = \pm \frac{\hbar}{2}.
\label{eq:chapter10-spin-half}
\]

Rotate the apparatus and the probabilities change, but the representation
structure remains. Spin is not a tiny classical arrow secretly pointing in an
ordinary direction. It is a representation of the rotation symmetry carried by
the quantum state.

\begin{phenom}{Spin One-Half Effect}
\label{ph:spin-one-half-effect}

\PhStatement A symmetry can have discrete representation content that no
continuous classical orientation can replace.

\PhOrigin Stern-Gerlach measurements revealed discrete spin outcomes for atoms
in magnetic-field gradients.

\PhObservation A spin one-half system measured along a chosen axis gives two
outcomes rather than a continuum of projected classical moments.

\PhConstraint The apparatus may not interpret the two outcomes as ignorance
about an ordinary vector when the rotation representation has different
structure.

\PhConsequence Invariance under rotation is carried by representation theory,
not by naive spatial imagery.

\end{phenom}

This is the representation-theory pleasure of the chapter. Physics is not
asking which cartoon picture survives. It is asking which structure is
invariant under admissible changes of basis.

\section{Yang--Mills and Confinement}
\label{se:yang-mills-and-confinement}

Yang-Mills theory generalizes gauge structure to nonabelian groups. The strong
interaction is organized by SU(3) color. Unlike electromagnetism, the gauge
bosons themselves carry the charge of the interaction. The field interacts with
itself.

Confinement is the operational shock. Isolated quarks are not observed. Try to
separate a quark pair, and the energy in the field grows until producing new
hadrons is cheaper than isolating a single colored object.

\begin{phenom}{Yang-Mills Confinement Effect}
\label{ph:yang-mills-confinement-effect}

\PhStatement A gauge charge may be physically real while never appearing as an
isolated admissible detector record.

\PhOrigin Nonabelian Yang-Mills theories underlie the strong interaction, where
color charge is confined inside hadrons.

\PhObservation Detectors see jets and hadrons, not free quarks, even though
quark structure explains scattering and spectroscopy.

\PhConstraint The measurement may not demand every internal gauge label to
appear as an isolated external object.

\PhConsequence Gauge invariance can make the observable record a composite
closure rather than a list of detachable constituents.

\end{phenom}

This is an important warning for the whole trilogy. Hidden structure is not
automatically illicit. It becomes illicit when it is used without paying the
measurement bill. Quarks pay their bill through hadron spectra, scattering, and
jets. They do not pay it by walking alone into a detector.

\section{Chirality and Parity Violation}
\label{se:chirality-and-parity-violation}

For a long time, physicists expected nature to respect mirror symmetry. If an
experiment was reflected left-to-right, the mirrored version should be equally
admissible. The weak interaction broke that expectation.

In the Wu experiment, cobalt-60 nuclei were aligned at low temperature, and the
directions of emitted beta particles were measured. The emission preferred one
direction relative to the nuclear spin. The mirror image was not equivalent.

\begin{phenom}{Chirality Effect}
\label{ph:chirality-effect}

\PhStatement A transformation that looks like an innocent relabeling can fail
operationally when the interaction distinguishes handedness.

\PhOrigin Parity violation in weak interactions was established experimentally
by Wu and collaborators using beta decay of aligned cobalt-60 nuclei.

\PhObservation Electrons were emitted preferentially opposite the nuclear spin,
showing that the weak interaction distinguishes left from right.

\PhConstraint The theory may not assume mirror symmetry as admissible when the
measurement record breaks it.

\PhConsequence Chirality is not decorative naming. For the weak interaction,
handedness is part of the physical record.

\end{phenom}

The chapter's answer is now complete. What remains invariant is not whatever
feels aesthetically symmetric. It is what the apparatus, conservation laws, and
interaction structure preserve. Some relabelings are gauges. Some are broken.
Some are forbidden by the record.

\begin{bridge}

The chapter's question was what remains invariant under admissible relabeling
or recompilation.

The answer is conserved and represented measurement content. Noether turns
symmetry into conservation. The photoelectric effect reveals quantized exchange.
Stern-Gerlach makes representation discrete. Yang-Mills confinement shows that
internal gauge labels can be real without appearing alone. The Wu experiment
shows that mirror relabeling is not always admissible.

The Lean boundary is explicit: this chapter does not claim the current code
formalizes Standard Model gauge groups. It uses the code's \lean{LOGICAL},
\lean{ComputerProgram}, \lean{MEASURED}, and \lean{COMPILED} interfaces as the
measurement grammar for future formalization.

The final physical question is thermodynamic. Once records are measured,
compiled, and conserved, why do they grow in one direction, and when are they
closed enough to support inference?

\end{bridge}

\begin{coda}{A Choreographed Performance Surviving Costume and Stage Changes}

The dance company performs a piece that was made twenty years ago. The original
cast has retired. The costumes have changed. The stage is smaller than the one
used at the premiere. The lighting designer has rebuilt the cues with modern
equipment.

Still, everyone says it is the same work.

Sameness does not mean that every visible detail survived. The dancers are not
the same people. The fabric is not the same fabric. The stage markings are not
the same tape. What survives is relational: who crosses whom, which gesture
answers which phrase, where weight transfers, how the ensemble turns, when the
silence lands.

The choreographer's notation is a gauge rule. It permits some substitutions
and forbids others. A taller dancer may take a shorter step if the spacing
remains correct. A costume may change color if it does not hide the line of the
arm. A lighting cue may move by a fraction of a second if the musical relation
is preserved. These are admissible relabelings.

Other changes break the work. Mirror the piece without permission and the
left-handed gestures become right-handed. Remove a counterweighting dancer and
a lift no longer carries the same risk. Change the rhythm under a turn and the
phrase loses its conservation. The performance has not been relabeled. It has
become another object.

This performs the chapter. Noether is the conservation of choreographic
relation under admissible change. Spin is the representation of turning that
cannot be replaced by a flat drawing of an arrow. Chirality is the left-handed
phrase that the mirror cannot preserve. Confinement is the ensemble figure that
exists only as a group; no single dancer can walk offstage carrying it alone.

Compilation appears in rehearsal. The director translates old notation into
living bodies. The result must be executable on this cast, this stage, this
night. A change is accepted only when the audience-facing work remains
invariant. The compiler has succeeded when the performance is new in its
materials and the same in its measured structure.

The curtain falls. The costumes go back to the rack. The dancers change. The
piece remains available because the admissible invariants were carried.

\end{coda}
